\documentclass[11pt,a4paper]{article}
\usepackage[utf8]{inputenc}
\usepackage[polish]{babel}
\usepackage[T1]{fontenc}
\usepackage{graphicx}
\usepackage{hyperref}
\usepackage{listings}
\usepackage{xcolor}
\usepackage{geometry}
\usepackage{float}
\usepackage{amsmath}

\geometry{margin=2.5cm}

\title{Issues}
\author{Fire Simulation System}
\date{\today}

\begin{document}

\maketitle
\newpage

\section{Wprowadzenie}

\subsection{Cel dokumentu}

Dokument zawiera opisy wszystkich błędów które nie udało się nam naprawić / poprawić.
\subsection{Zakres dokumentu}

\section{Zatrzymywanie symulacji}
Dalej pojawiają się błędy przy zatrzymywaniu symulacji / supportu, w wątkach 
odpowiedzialnych za odbieranie wiadomości z RabbitMQ. Wymusza to restart symulatora i supportu przy każdej nowej symulacji. 

\section{Niedokończone funkcjonalności}
Brak możliwości modyfikacji ustawień symulatora w wizualizacji. Przygotowany formularz nie jest 
zintegrowany w pełni z symulatorem oraz backendem. 

\section{Brak możliwości dynamicznego zmieniania typu rekomendacji}
Zmiana z MCTS na LLM wymaga restartu systemu z zmianą zmiennych środowiskowych 

\section{Problemy wydajnościowe}
MCTS dla symulacji 25x25 (625 sektorów) generuje różne, często sprzeczne rekomendacje. Wynika to z ograniczeń na czas (1s domyślnie), oraz z ogromu sektorów do przetworzenia. Brygady potrafią otrzymać więc różne rekomendowane akcje w danym momencie. 

\section{Kod}
Kod odpowiedzialny za wizualizację wymaga refactoringu. Poruszanie się po nim jest praktycznie niemożliwe, zalecamy próbę przepisania go od zera. 

\section{Logi}
Nie udało się nam opracować stabilnego modelu logowania, logi w trybie DEBUG są zbyt szczegółowe, INFO potrafią nie wyjaśniać błędów symulacji. 

\section{RabbitMQ}
Zbyt duża liczba kolejek - wymagane uproszczenie dla ilości wszystkich tworzonych kolejek, aby proces komunikacji był bardziej przejrzysty i prostszy do debugu. 

\section{LLM}
Rekomendacje LLM-owe, generowanie komunikatów językiem naturalnym dla brygad wymaga odpytywania API które trwa zbyt długo i powoduje że symulacja nie jest dynamiczna. Warto rozważyć postawienie lokalnego LLM-a - np. przetestować zastosowanie https://ollama.com/ dla tego projektu. 
\end{document}
